% !TeX program = xelatex
% 全国大学生数学建模竞赛论文结构模板
% 说明：本模板只保留优秀论文中高频出现的行文框架，不预置具体建模内容。
% 建议在 Overleaf 中选择 XeLaTeX 编译器。

\documentclass[withoutpreface,bwprint]{cumcmthesis}

% -------------------- 标准竞赛排版 --------------------
\usepackage{cumcm2026}
\usepackage{amsmath,amssymb,bm}
\usepackage{graphicx}
\usepackage{booktabs}
\usepackage{array}
\usepackage{float}
\usepackage{enumitem}
\usepackage{hyperref}
\usepackage{gbt7714}
\usepackage[section]{placeins}
\hypersetup{hidelinks}
\title{中药材热风烘干过程的传热传质建模与烘干时长预测}

% 交叉引用：将“表”和表号整体设置为可点击链接。
\newcommand{\tabref}[1]{\hyperref[#1]{表~\ref*{#1}}}
\newcommand{\figref}[1]{\hyperref[#1]{图~\ref*{#1}}}

% -------------------- 基础版式 --------------------
\setlength{\parindent}{2em}
\setlength{\parskip}{0pt}

% 若需要统一图表文件夹，可直接使用 figures/ 下的文件
\graphicspath{{figures/}}
\newcommand{\paperfigure}[4][0.92]{%
  \begin{figure}[!htbp]\centering
  \includegraphics[width=#1\linewidth,height=0.34\textheight,keepaspectratio]{#2}
  \caption{#3}\label{#4}
  \end{figure}}
\newcommand{\unitC}{\,{}^\circ\mathrm{C}}
% 图4、图10横向放大；图6同时解除原来的高度限制。
\newcommand{\largepaperfigure}[4][1.10]{%
  \begin{figure}[!htbp]\centering
  \makebox[\linewidth][c]{\includegraphics[width=#1\linewidth,height=0.55\textheight,keepaspectratio]{#2}}
  \caption{#3}\label{#4}
  \end{figure}}
\setlength{\emergencystretch}{2em}
\renewcommand{\topfraction}{0.9}
\renewcommand{\bottomfraction}{0.8}
\renewcommand{\textfraction}{0.08}
\renewcommand{\floatpagefraction}{0.75}
\raggedbottom

\begin{document}

% 国赛电子版论文第一页直接从摘要开始。
% 来源：标题摘要关键词_Loop_V3.docx。
\begin{center}
{\zihao{3}\heiti\bfseries 中药材热风烘干过程的传热传质建模与烘干时长预测}
\end{center}
\vspace{0.5em}
\begin{center}{\heiti\zihao{4} 摘\quad 要}\end{center}

本文通过有限体积法（FVM）进行空间离散，并采用变阶、自适应步长的后向差分公式（BDF）进行时间积分，对圆柱形中药材的热风烘干过程进行分析。逐步考虑密度、比热容、导热系数和水分扩散系数的变化及径向收缩，着重分析中部截面的温度和含水率分布。

\par\vspace{\baselineskip}
{\heiti\bfseries 针对问题一，}在药材尺寸及密度、比热容、导热系数保持不变、水分扩散系数随含水率变化的条件下，分别建立温度场和含水率场模型。使用分段三次Hermite插值（PCHIP）\cite{fritsch1980monotone}，得到环境温度和水分浓度关于时间的连续函数，并采用对流换热和对流传质边界，求解预热阶段的场量分布。数值结果表明，在$1800\,\mathrm{s}$时，中部截面中心与表面温度分别为$33.5758\unitC$和$36.7858\unitC$，含水率分别为$2.5500\,\mathrm{kg/kg}$和$1.5103\,\mathrm{kg/kg}$。

\par\vspace{\baselineskip}
{\heiti\bfseries 针对问题二，}在药材尺寸保持不变的条件下，根据当前温度和含水率更新密度、比热容、导热系数和水分扩散系数，联立求解温度场和含水率场。数值结果表明，在$3\,\mathrm{h}$时，中心与表面温度分别为$49.8506\unitC$和$49.9674\unitC$，仅相差$0.1168\unitC$；含水率分别为$1.7662\,\mathrm{kg/kg}$和$1.0081\,\mathrm{kg/kg}$，相差$0.7581\,\mathrm{kg/kg}$。药材的热平衡明显早于水分平衡，后续烘干主要受内部水分向外扩散控制。

\par\vspace{\baselineskip}
{\heiti\bfseries 针对问题三，}将附件环境数据末20组成对观测的平均值作为长时烘干的环境条件，以全域最大含水率首次降至$0.15\,\mathrm{kg/kg}$的时刻作为干燥终点。计算得到烘干总时间为$57.4354\,\mathrm{h}$；全域扫描结果显示，几何中心最后达到干燥阈值。

\par\vspace{\baselineskip}
{\heiti\bfseries 针对问题四，}在径向收缩的动态区域上重新求解温度场和含水率场，并沿用问题三的环境延拓条件和干燥判据。计算得到烘干时间为$51.0583\,\mathrm{h}$，结束时药材半径为$1.2000\,\mathrm{cm}$。与问题三相比，烘干时长减少$6.3772\,\mathrm{h}$，即$11.10\%$；该时间差反映了半径变化及密度、比热容、导热系数和水分扩散系数变化的共同作用。

\par\vspace{\baselineskip}
\noindent\textbf{关键词：}中药材烘干；传热；传质；有限体积法；径向收缩

\clearpage

% ============================================================
% 一、问题重述
% ============================================================
\section{问题重述}

\subsection{问题背景}
干燥是中药材加工中的重要工序。热风干燥因设备简单、操作方便而应用广泛，但其干燥效率和成品品质易受工艺条件影响\cite{LiWang2024MedicinalDrying,Li2021GastrodiaDrying}。

热风干燥过程中，药材与热空气持续进行热量和水分交换。烘干工艺可分为预热平衡和恒温干燥两个阶段；从干燥动力学看，过程通常表现为升温、恒速和降速三个阶段\cite{Yan2026DryingHeatMassTransfer}。药材内部温度和水分浓度随时间与空间变化，是描述完整烘干过程的重要依据。

随着水分持续流失，药材可能发生收缩并引起几何尺寸变化，进而影响温度和水分分布的预测\cite{MayorSereno2004Shrinkage}。烘房条件、药材参数及尺寸变化共同影响实际烘干过程。

传统依赖反复试验调节工艺参数的方法成本高、周期长。有必要借助数理分析与数值仿真，揭示药材的烘干规律，为过程预测与烘干终点判定提供依据。

\subsection{问题提出}
研究对象为近似圆柱形中药材。根据烘房温湿环境、药材参数及尺寸变化等已知信息，需要解决以下四个问题。

\noindent\textbf{问题一：}建立预热平衡阶段的数学模型，描述药材内部温度和水分浓度随时间、径向位置的变化规律。

\noindent\textbf{问题二：}考虑预热平衡向恒温干燥过渡时的参数变化，将模型拓展至完整烘干过程，并分析初始 $3~\mathrm{h}$ 内温度和水分浓度的演化特征。

\noindent\textbf{问题三：}以药材各处水分浓度均低于 $0.15~\mathrm{kg/kg}$ 作为烘干完成判据，根据水分浓度演化结果确定所需时长。

\noindent\textbf{问题四：}进一步考虑水分流失引起的药材尺寸收缩，并依据半径变化修正模型，重新研究水分迁移过程及烘干时长。

% 来源：问题分析_Word版.docx。
\section{问题分析}
\subsection{问题一的分析}
问题一是后续各问的基础，其核心是描述预热平衡阶段药材内部的热量与水分传递过程。仅考虑径向变化能够描述径向中心与表面的差异，但无法反映端面交换及轴向梯度。因此，将药材视为二维轴对称区域，同时保留径向和轴向的传热传质。预热阶段药材的密度、比热容、导热系数按题目附录2取定值，水分扩散系数随局部含水率变化；烘房温度和水分浓度由附件1确定。

该问的求解重点是在变化的环境条件下计算温度场和含水率场，并获得指定时刻、位置的结果。本文采用有限体积法进行空间离散，利用自适应BDF方法进行时间积分，由此建立后续问题共用的计算框架。

计算得到整个二维区域内的场量后，选取轴向中截面$z=0$进行展示。温度场和含水率场关于该截面对称，该截面距离两个端面最远，受端面换热和水分交换的直接影响相对较弱，能够反映药材内部的变化。因此，四个问题均提取$T(r,0,t)$和$C(r,0,t)$作为代表性输出。

\subsection{问题二的分析}
在问题一的基础上，采用题目附录3给出的物性关系，将研究范围扩展到初始$3\,\mathrm{h}$烘干阶段。药材密度、比热容和导热系数随局部含水率变化，水分扩散系数同时受到温度和含水率的影响。含水率下降改变传热参数，温度变化又影响扩散速度，使两个物理场相互关联。因此，在积分过程中根据当前状态更新相应系数，同步求解温度场和含水率场。

\subsection{问题三的分析}
问题三要求确定药材达到烘干标准所需的总时间，其传热传质机理与问题二相同，但需要将计算持续推进至药材内部各处均满足烘干要求。为此，持续监测全域最大含水率，并以其首次降至阈值的时刻作为总烘干时间。

附件1仅给出了前$4\,\mathrm{h}$的烘房环境数据，而题目指出实际烘干通常持续2至3天，长期求解必须明确处理数据截止后的环境条件。本文将末20组成对观测的温度和水分浓度均值作为后续环境条件，并通过改变平均窗口检查烘干时间对该选择的敏感程度。

\subsection{问题四的分析}
问题四进一步考虑药材尺寸变化，需要在更新本问密度、比热容、导热系数和水分扩散系数的同时，根据附件2插值得到半径函数$R(t)$。假设药材长度保持不变，仅在径向均匀收缩，将当前径向位置按半径归一化，使节点始终对应药材内部的同一相对位置。计算中只需更新节点实际位置、控制体体积和界面面积，无需重新生成网格连接关系。

烘干终点继续采用全域最大含水率首次达到规定阈值的判据。最后，将本问的烘干时间与问题三比较，分析径向收缩和物性变化对烘干过程的综合影响。

\section{模型假设}
\begin{enumerate}[label=\arabic*.]
\item 药材为规则圆柱体，内部介质连续、均匀且各向同性，场量仅随径向位置、轴向位置及时间变化。
\item 初始温度与干基含水率在药材内部均匀分布；周向环境一致，两个端面的换热和传质条件相同。
\item 药材表面满足对流换热和对流传质边界条件，对流换热系数与对流传质系数保持不变。
\item 内部热量通过热传导传递，水分通过等效扩散迁移；忽略内部对流、化学反应和蒸发潜热对温度方程的直接贡献。
\item 问题一至问题三保持药材尺寸不变；问题四长度保持不变，仅考虑径向均匀收缩，场量在随骨架移动的材料坐标下演化。
\item 环境数据范围内采用PCHIP插值；问题三和问题四在数据截止后采用末20组成对观测均值延拓环境条件。
\end{enumerate}

\section{符号说明}
本文以$C$表示药材干基含水率，即题目中的水分浓度。正文中的温度$T$以摄氏度表示，扩散系数经验式中的绝对温度统一写为$T+273.15$。主要符号及单位见\tabref{tab:symbols}。
\begin{table}[!htbp]\centering\small
\caption{主要符号及其含义与单位}\label{tab:symbols}
\renewcommand{\arraystretch}{1.12}
\begin{tabularx}{0.98\textwidth}{@{}cLc@{}}
\toprule
符号 & 含义 & 单位\\\midrule
$t$ & 自烘干开始计的时间 & $\mathrm{s}$\\
$r,z$ & 径向坐标、以中截面为原点的轴向坐标 & $\mathrm{m}$\\
$R(t),R_0$ & 当前半径、初始半径 & $\mathrm{m}$\\
$L,H$ & 圆柱长度、半长（$H=L/2$） & $\mathrm{m}$\\
$T,T_\infty$ & 药材温度、烘房温度 & $ {}^\circ\mathrm{C}$\\
$C,C_\infty$ & 药材干基含水率、烘房水分浓度 & $\mathrm{kg/kg}$\\
$\rho$ & 药材密度 & $\mathrm{kg/m^3}$\\
$c_p$ & 药材比热容 & $\mathrm{J/(kg\cdot K)}$\\
$k$ & 导热系数 & $\mathrm{W/(m\cdot K)}$\\
$D$ & 有效水分扩散系数 & $\mathrm{m^2/s}$\\
$h$ & 对流换热系数 & $\mathrm{W/(m^2\cdot K)}$\\
$h_m$ & 对流传质系数 & $\mathrm{m/s}$\\
$N_r,N_z$ & 径向、轴向网格区间数 & 1\\
$\lambda$ & 网格拉伸指数 & 1\\
$V_P,A_f,d_{PN}$ & 控制体体积、界面面积、相邻节点间距 & $\mathrm{m^3},\mathrm{m^2},\mathrm{m}$\\
$\xi,s_R$ & 归一化径向坐标、半径缩放比例 & 1\\
$m$ & 环境延拓所用末段观测点数 & 1\\
$t_{\mathrm{dry}}$ & 首次达到全域干燥阈值的时刻 & $\mathrm{s}$或$\mathrm{h}$\\
\bottomrule
\end{tabularx}
\end{table}

% 来源：第六部分_图文整合框架_Loop_V4.docx（含全部7条批注及汇总核对清单）。
% 实际章节编号由主文件顺序自动生成为第五部分。
\section{模型的建立与求解}
\begingroup
% V4批注：压缩公式与表格周围留白，以及成组公式行距。
\setlength{\abovedisplayskip}{6pt plus 1pt minus 1pt}
\setlength{\belowdisplayskip}{6pt plus 1pt minus 1pt}
\setlength{\abovedisplayshortskip}{3pt}
\setlength{\belowdisplayshortskip}{4pt}
\setlength{\jot}{2pt}
\setlength{\intextsep}{9pt plus 2pt minus 2pt}
\setlength{\textfloatsep}{11pt plus 2pt minus 2pt}
\setlength{\floatsep}{8pt plus 2pt minus 2pt}
\subsection{数值求解框架}
四个问题对应的密度、比热容、导热系数、水分扩散系数、几何条件与烘干时长不同，但均可归结为药材内部的非稳态热传导与水分扩散问题。本文采用二维轴对称有限体积法进行空间离散，采用变阶、自适应步长的BDF方法进行时间积分。问题一分别求解两个物理场；问题二至问题四根据局部温度和含水率更新密度、比热容、导热系数和水分扩散系数，联立求解。

\subsubsection{计算区域的化简}
如\figref{fig:geometry}所示，将药材近似为半径$R$、长度$L$的均质圆柱体，以轴向中截面为$z=0$。圆柱初始半径$R_0=0.02\,\mathrm{m}$，长度$L=0.25\,\mathrm{m}$，半长$H=L/2=0.125\,\mathrm{m}$。在周向环境一致、两个端面条件相同的假设下，场量与周向坐标$\theta$无关，即
\begin{equation}
\frac{\partial T}{\partial\theta}=\frac{\partial C}{\partial\theta}=0.
\end{equation}
\noindent 利用中截面对称性，只需在半个子午区域
\begin{equation}\label{eq:domain}
\Omega=\{(r,z)\mid 0\le r\le R,\ 0\le z\le H\}
\end{equation}
\noindent 内求解。$r=0$和$z=0$为对称边界，$r=R$和$z=H$分别对应侧壁与端面。后两者均与烘房交换热量和水分。场量$u(r,z,t)$的扩散算子为
\begin{equation}\label{eq:operator}
\nabla\cdot(a\nabla u)=
\frac1r\frac{\partial}{\partial r}\left(ra\frac{\partial u}{\partial r}\right)
+\frac{\partial}{\partial z}\left(a\frac{\partial u}{\partial z}\right).
\end{equation}
\paperfigure{cylinder_geometry_symbols_dot_only.png}{圆柱体几何参数}{fig:geometry}

\subsubsection{初值与边界条件}
初始温度和含水率均匀分布：
\begin{equation}\label{eq:initial}
T(r,z,0)=T_0=28\unitC,\qquad
C(r,z,0)=C_0=2.55\,\mathrm{kg/kg}.
\end{equation}
\noindent 中心轴和中截面的法向通量均为零，因而有
\begin{equation}\label{eq:symmetry}
\left.\frac{\partial T}{\partial r}\right|_{r=0}
=\left.\frac{\partial C}{\partial r}\right|_{r=0}=0,\qquad
\left.\frac{\partial T}{\partial z}\right|_{z=0}
=\left.\frac{\partial C}{\partial z}\right|_{z=0}=0.
\end{equation}
\noindent 侧壁与端面的热交换采用牛顿冷却定律，水分交换采用对流传质关系。以$\boldsymbol n$表示外法向，两类外边界统一写为
\begin{equation}\label{eq:robin}
-k\frac{\partial T}{\partial n}=h(T-T_\infty(t)),\qquad
-D\frac{\partial C}{\partial n}=h_m(C-C_\infty(t)).
\end{equation}
\noindent 其中，侧壁上$\partial/\partial n=\partial/\partial r$，端面上$\partial/\partial n=\partial/\partial z$；题目给定
\begin{equation}
h=25\,\mathrm{W/(m^2\cdot K)},\qquad
h_m=8\times10^{-7}\,\mathrm{m/s}.
\end{equation}
\noindent 对附件1的离散环境数据采用分段三次Hermite插值（Piecewise Cubic Hermite Interpolating Polynomial，PCHIP）\cite{fritsch1980monotone}，构造连续边界函数$T_\infty(t)$和$C_\infty(t)$，避免在单调数据区间引入额外过冲。

\subsubsection{有限体积离散与时间积分}
将两类场量统一写为
\begin{equation}\label{eq:unified}
b_u(T,C)\frac{\partial u}{\partial t}
=\nabla\cdot\bigl(a_u(T,C)\nabla u\bigr).
\end{equation}
温度场对应$u=T$、$b_T=\rho c_p$、$a_T=k$；含水率场对应$u=C$、$b_C=1$、$a_C=D$。

考虑到外表面附近梯度较大，采用向侧壁和端面加密的节点分布：
\begin{equation}\label{eq:grid}
\begin{aligned}
&s_i=\frac{i}{N_r},\quad
r_i=R\left[1-(1-s_i)^\lambda\right],&&i=0,\ldots,N_r,\\
&q_j=\frac{j}{N_z},\quad
z_j=H\left[1-(1-q_j)^\lambda\right],&&j=0,\ldots,N_z.
\end{aligned}
\end{equation}
其中$N_r,N_z$为网格区间数，$\lambda=1$对应均匀网格，$\lambda>1$对应近壁加密；正式计算取$\lambda=2$。

以节点$P=(r_i,z_j)$为中心构造控制体，内部公共界面取相邻节点的中点：
\begin{equation}
r_{i+\frac12}=\frac{r_i+r_{i+1}}2,\qquad
z_{j+\frac12}=\frac{z_j+z_{j+1}}2.
\end{equation}
边界节点的外侧界面与实际边界重合，即
$r_{-\frac12}=0$、$r_{N_r+\frac12}=R$、$z_{-\frac12}=0$及$z_{N_z+\frac12}=H$。
令$\Delta z_j=z_{j+\frac12}-z_{j-\frac12}$，控制体绕轴旋转后的体积与界面面积为
\begin{equation}\label{eq:volume}
\begin{aligned}
V_P&=\pi\left(r_{i+\frac12}^2-r_{i-\frac12}^2\right)\Delta z_j,\\
A^r_{i+\frac12,j}&=2\pi r_{i+\frac12}\Delta z_j,\\
A^z_{i,j+\frac12}&=\pi\left(r_{i+\frac12}^2-r_{i-\frac12}^2\right).
\end{aligned}
\end{equation}

对式\eqref{eq:unified}在控制体上积分，应用高斯散度定理，有
\begin{equation}
\int_{V_P}b_u\frac{\partial u}{\partial t}\,\mathrm dV
=\int_{\partial V_P}a_u\nabla u\cdot\boldsymbol n\,\mathrm dS.
\end{equation}
\figref{fig:fvm}以传热为例说明控制体之间的通量交换，图中的$Q$表示通过界面的热流率。用节点值近似控制体内的场量与时间项系数，左端近似为$b_{u,P}V_P\,\mathrm du_P/\mathrm dt$。
\paperfigure{fvm_multiscale_macro_micro.png}{网格和通量交换示意图}{fig:fvm}

对相邻控制体$P$与$N$的公共界面$f$，采用中心差分和扩散系数的算术平均：
\begin{equation}
\left.\frac{\partial u}{\partial n}\right|_f
\approx\frac{u_N-u_P}{d_{PN}},\qquad
a_{u,f}=\frac{a_{u,P}+a_{u,N}}2.
\end{equation}
因此，由相邻节点$N$流入控制体$P$的通量为
\begin{equation}
F_{PN}=\frac{a_{u,f}A_f}{d_{PN}}(u_N-u_P).
\end{equation}
对各界面求和，得到统一的半离散方程
\begin{equation}\label{eq:fvm}
b_{u,P}V_P\frac{\mathrm du_P}{\mathrm dt}
=\sum_{N\in\mathcal N_P}
\frac{a_{u,f}A_f}{d_{PN}}(u_N-u_P)+\Phi_{b,P}.
\end{equation}
其中$\mathcal N_P$为相邻节点集合。对称边界的交换通量为零；对第三类边界，
\begin{equation}
\Phi_{b,P}=\sum_{b\subset\partial V_P\cap\partial\Omega_{\rm ext}}
\beta A_{b,P}(u_\infty-u_P).
\end{equation}
温度场取$\beta=h$，含水率场取$\beta=h_m$。位于侧壁与端面交接处的控制体同时累加两部分边界通量。该积分形式也避免了在中心轴上直接计算式\eqref{eq:operator}中的$1/r$项。

将所有节点的温度与含水率排列为状态向量
\begin{equation}\label{eq:ode}
\boldsymbol y(t)=
\begin{bmatrix}\boldsymbol T(t)\\\boldsymbol C(t)\end{bmatrix},
\qquad
\frac{\mathrm d\boldsymbol y}{\mathrm dt}
=\boldsymbol f(t,\boldsymbol y).
\end{equation}
由于传热与水分扩散时间尺度差异明显，且密度、比热容、导热系数和水分扩散系数随状态变化，所得方程组具有刚性。采用变阶、自适应步长BDF方法，并利用稀疏雅可比矩阵推进积分；相对与绝对误差容限分别取$10^{-7}$和$10^{-9}$。

\FloatBarrier
\subsection{问题一：预热平衡阶段}
\subsubsection{温度场和含水率场模型}
预热阶段的密度、比热容和导热系数分别为
\begin{equation}
\rho=820\,\mathrm{kg/m^3},\qquad
c_p=2600\,\mathrm{J/(kg\cdot K)},\qquad
k=0.36\,\mathrm{W/(m\cdot K)}.
\end{equation}
温度场满足
\begin{equation}\label{eq:q1-temperature}
\rho c_p\frac{\partial T}{\partial t}
=\frac1r\frac{\partial}{\partial r}
\left(rk\frac{\partial T}{\partial r}\right)
+\frac{\partial}{\partial z}\left(k\frac{\partial T}{\partial z}\right).
\end{equation}
含水率场采用有效扩散模型：
\begin{equation}\label{eq:q1-moisture}
\frac{\partial C}{\partial t}
=\frac1r\frac{\partial}{\partial r}
\left(rD(C)\frac{\partial C}{\partial r}\right)
+\frac{\partial}{\partial z}\left(D(C)\frac{\partial C}{\partial z}\right),
\end{equation}
其中
\begin{equation}\label{eq:q1-D}
D(C)=7\times10^{-9}\exp\left(-\frac{0.89}{C}\right)\quad(\mathrm{m^2/s}).
\end{equation}
初始与边界条件见式\eqref{eq:initial}至式\eqref{eq:robin}。温度方程不含含水率变量，水分方程不含温度变量，因此两者可以解耦，分别积分。

\paperfigure{fig1_chamber_temp_humidity.png}{烘房环境观测数据}{fig:environment}
\figref{fig:environment}中散点为观测值，曲线为插值结果。环境温度和水分浓度在前期上升，随后趋于稳定并伴随小幅波动。

\subsubsection{模型求解与结果分析}
将$b_T=\rho c_p$、$a_T=k$和$b_C=1$、$a_C=D(C)$代入式\eqref{eq:fvm}。采用$N_r=N_z=120$、$\lambda=2$的网格，从$t=0$积分至$t=1800\,\mathrm{s}$。按照题目规定，提取中截面上$0$、$0.5$、$1.0$、$1.5$、$2.0\,\mathrm{cm}$处的结果，见\tabref{tab:q1-temperature}与\tabref{tab:q1-moisture}。
\begin{table}[!htbp]\centering\small
\caption{预热阶段中部截面的温度，单位：$ {}^\circ\mathrm{C}$}\label{tab:q1-temperature}
\renewcommand{\arraystretch}{1.15}
\begin{tabular*}{0.96\textwidth}{@{\extracolsep{\fill}}cccccc}
\toprule
 & \multicolumn{5}{c}{到药材轴线的距离/cm}\\
\cmidrule(lr){2-6}
时间/s & 0 & 0.5 & 1.0 & 1.5 & 2.0\\\midrule
100.0000 & 28.0000 & 28.0003 & 28.0039 & 28.0319 & 28.1795\\
300.0000 & 28.0406 & 28.0634 & 28.1513 & 28.3672 & 28.8486\\
600.0000 & 28.4534 & 28.5361 & 28.8039 & 29.3155 & 30.1640\\
900.0000 & 29.3243 & 29.4583 & 29.8757 & 30.6164 & 31.7295\\
1200.0000 & 30.5429 & 30.7099 & 31.2225 & 32.1130 & 33.4294\\
1500.0000 & 31.9960 & 32.1870 & 32.7667 & 33.7476 & 35.1211\\
1800.0000 & 33.5758 & 33.7724 & 34.3642 & 35.3621 & 36.7858\\
\bottomrule
\end{tabular*}
\end{table}

\begin{table}[!htbp]\centering\small
\caption{预热阶段中部截面的干基含水率，单位：$\mathrm{kg/kg}$}\label{tab:q1-moisture}
\renewcommand{\arraystretch}{1.15}
\begin{tabular*}{0.96\textwidth}{@{\extracolsep{\fill}}cccccc}
\toprule
 & \multicolumn{5}{c}{到药材轴线的距离/cm}\\
\cmidrule(lr){2-6}
时间/s & 0 & 0.5 & 1.0 & 1.5 & 2.0\\\midrule
100.0000 & 2.5500 & 2.5500 & 2.5500 & 2.5500 & 2.2472\\
300.0000 & 2.5500 & 2.5500 & 2.5500 & 2.5492 & 2.0509\\
600.0000 & 2.5500 & 2.5500 & 2.5500 & 2.5353 & 1.8770\\
900.0000 & 2.5500 & 2.5500 & 2.5497 & 2.5046 & 1.7548\\
1200.0000 & 2.5500 & 2.5500 & 2.5482 & 2.4646 & 1.6586\\
1500.0000 & 2.5500 & 2.5499 & 2.5445 & 2.4206 & 1.5788\\
1800.0000 & 2.5500 & 2.5497 & 2.5383 & 2.3755 & 1.5103\\
\bottomrule
\end{tabular*}
\end{table}


在$1800\,\mathrm{s}$时，中心温度为$33.5758\unitC$，表面温度为$36.7858\unitC$；中心含水率仍为$2.5500\,\mathrm{kg/kg}$，表面已降至$1.5103\,\mathrm{kg/kg}$。热量传递已影响整个中截面，而明显失水主要集中于外侧区域。

\largepaperfigure[1.18]{q1_combined_pizza.png}{预热阶段中部截面的含水率与温度变化（左：含水率；右：温度）}{fig:q1-fields}
\figref{fig:q1-fields}将时间沿圆周顺时针展开，距圆心的距离表示径向位置$r$，颜色表示含水率或温度。表层首先与环境交换水分，失水影响逐步向内扩展，中心保持较高含水率；表面也最先升温，随后热量向内部传导。

\subsubsection{收敛性验证}
保持密度、比热容、导热系数、水分扩散关系、环境边界和积分容限不变，比较每个方向分别含60、80、100、120、140和160个区间的网格。以$160\times160$网格的输出为参考，定义
\begin{equation}\label{eq:error}
\begin{aligned}
E_T(N)&=\max_{t_k,r_l}|T_N(r_l,0,t_k)-T_{160}(r_l,0,t_k)|,\\
E_C(N)&=\max_{t_k,r_l}|C_N(r_l,0,t_k)-C_{160}(r_l,0,t_k)|.
\end{aligned}
\end{equation}
其中$t_k$为输出时刻，$r_l$为$0$至$2\,\mathrm{cm}$内的指定位置。比较结果见\tabref{tab:q1-grid}。
\begin{table}[!htbp]\centering\small
\caption{问题一各网格相对参考网格的最大差异}\label{tab:q1-grid}
\renewcommand{\arraystretch}{1.15}
\begin{tabular*}{0.96\textwidth}{@{\extracolsep{\fill}}ccc}
\toprule
网格区间数 & 温度最大差/$ {}^\circ\mathrm{C}$ & 含水率最大差/(kg/kg)\\\midrule
$ 60\times 60$ & 0.0002 & 0.0004\\
$ 80\times 80$ & 0.0001 & 0.0002\\
$ 100\times 100$ & 0.0001 & 0.0001\\
$ 120\times 120$ & 0.0001 & 0.0001\\
$ 140\times 140$ & 0.0001 & 0.0001\\
$ 160\times 160$ & 0.0000 & 0.0000\\
\bottomrule
\end{tabular*}
\par\smallskip{\footnotesize 注：比较现有四位小数导出值，参考网格为$160\times160$；表中差异受导出精度限制。}
\end{table}


$120\times120$与$160\times160$网格的两项最大差均处于$10^{-4}$量级，继续加密对输出结果的影响较小，在当前输出精度下可以认为已经收敛。由于现有导出值保留四位小数，该比较不等同于未舍入解的严格误差上界。

均匀网格与近壁加密网格的进一步比较表明，以各自$160\times160$结果为参考，$60\times60$均匀网格的含水率最大差为$2.9\times10^{-3}\,\mathrm{kg/kg}$，相同规模的近壁加密网格为$4.0\times10^{-4}\,\mathrm{kg/kg}$；增至$120\times120$时，两者分别为$4.0\times10^{-4}$和$1.0\times10^{-4}\,\mathrm{kg/kg}$。最大差异集中在预热初期的表层，说明加密高梯度区域能提高相同规模网格的输出精度。

\FloatBarrier
\subsection{问题二：变物性温湿耦合模型}
\subsubsection{模型建立}
保持药材尺寸不变，按题目附录3更新密度、比热容、导热系数和水分扩散系数：
\begin{equation}\label{eq:q2-properties}
\begin{aligned}
\rho(C)&=650+128C,\\
c_p(C)&=1450+2736\frac{C}{C+1},\\
k(C)&=0.21+0.38\frac{C}{C+1},\\
D(T,C)&=2.4\times10^{-3}
\exp\left(-\frac{0.45}{C}\right)
\exp\left(-\frac{3850}{T+273.15}\right).
\end{aligned}
\end{equation}
上述各量的单位见\tabref{tab:symbols}。温度场和含水率场满足
\begin{equation}\label{eq:q2-pde}
\begin{aligned}
\rho(C)c_p(C)\frac{\partial T}{\partial t}
&=\frac1r\frac{\partial}{\partial r}
\left(rk(C)\frac{\partial T}{\partial r}\right)
+\frac{\partial}{\partial z}\left(k(C)\frac{\partial T}{\partial z}\right),\\
\frac{\partial C}{\partial t}
&=\frac1r\frac{\partial}{\partial r}
\left(rD(T,C)\frac{\partial C}{\partial r}\right)
+\frac{\partial}{\partial z}\left(D(T,C)\frac{\partial C}{\partial z}\right).
\end{aligned}
\end{equation}
含水率通过密度、比热容和导热系数影响温度场，温度与含水率又共同决定扩散系数，因此需要联立求解。

\subsubsection{模型求解与结果分析}
两个场的时间项系数与扩散通量系数分别为
\begin{equation}
b_T=\rho(C)c_p(C),\quad a_T=k(C),\qquad
b_C=1,\quad a_C=D(T,C).
\end{equation}
采用$N_r=N_z=120$、$\lambda=2$，两个方向各有121个节点。在计算方程右端与雅可比矩阵时，按当前状态更新密度、比热容、导热系数、水分扩散系数及界面通量，以BDF方法在$0\le t\le10800\,\mathrm{s}$内积分，误差容限沿用前述设置。

每隔$0.5\,\mathrm{h}$、径向每隔$0.5\,\mathrm{cm}$的结果见\tabref{tab:q2-temperature}与\tabref{tab:q2-moisture}。
\begin{table}[!htbp]\centering\small
\caption{初始3小时中部截面的温度，单位：$ {}^\circ\mathrm{C}$}\label{tab:q2-temperature}
\renewcommand{\arraystretch}{1.15}
\begin{tabular*}{0.96\textwidth}{@{\extracolsep{\fill}}cccccc}
\toprule
 & \multicolumn{5}{c}{到药材轴线的距离/cm}\\
\cmidrule(lr){2-6}
时间/h & 0 & 0.5 & 1.0 & 1.5 & 2.0\\\midrule
0.5000 & 32.1896 & 32.3824 & 32.9660 & 33.9608 & 35.4132\\
1.0000 & 40.3822 & 40.5542 & 41.0609 & 41.8783 & 43.0000\\
1.5000 & 45.8476 & 45.9357 & 46.1938 & 46.6059 & 47.1401\\
2.0000 & 48.4518 & 48.4896 & 48.5998 & 48.7747 & 49.0042\\
2.5000 & 49.4687 & 49.4809 & 49.5152 & 49.5658 & 49.6627\\
3.0000 & 49.8506 & 49.8565 & 49.8757 & 49.9110 & 49.9674\\
\bottomrule
\end{tabular*}
\end{table}

\begin{table}[!htbp]\centering\small
\caption{初始3小时中部截面的干基含水率，单位：$\mathrm{kg/kg}$}\label{tab:q2-moisture}
\renewcommand{\arraystretch}{1.15}
\begin{tabular*}{0.96\textwidth}{@{\extracolsep{\fill}}cccccc}
\toprule
 & \multicolumn{5}{c}{到药材轴线的距离/cm}\\
\cmidrule(lr){2-6}
时间/h & 0 & 0.5 & 1.0 & 1.5 & 2.0\\\midrule
0.5000 & 2.5499 & 2.5489 & 2.5255 & 2.3257 & 1.6487\\
1.0000 & 2.5257 & 2.4948 & 2.3578 & 2.0230 & 1.4711\\
1.5000 & 2.3860 & 2.3256 & 2.1344 & 1.8020 & 1.3476\\
2.0000 & 2.1708 & 2.1084 & 1.9236 & 1.6259 & 1.2311\\
2.5000 & 1.9566 & 1.9006 & 1.7360 & 1.4720 & 1.1167\\
3.0000 & 1.7662 & 1.7165 & 1.5702 & 1.3333 & 1.0081\\
\bottomrule
\end{tabular*}
\end{table}

在$3\,\mathrm{h}$时，中心与表面温度之差仅为$0.1168\unitC$，但含水率之差仍为$0.7581\,\mathrm{kg/kg}$，说明温度趋于均匀的过程明显早于水分均衡。

\paperfigure[0.82]{result2_moisture_lines_0p3cm.png}{中部截面不同径向位置的含水率随时间变化}{fig:q2-moisture}
\figref{fig:q2-moisture}展示中截面$z=0$上的含水率时序曲线。表层含水率较早下降，中心附近下降较晚，至$3\,\mathrm{h}$各位置之间仍存在明显差异。这一结果表明内部水分迁移仍需持续进行，为问题三延长积分时间提供依据。

\subsubsection{收敛性验证}
在温湿耦合模型下，沿用式\eqref{eq:error}比较不同网格的结果，见\tabref{tab:q2-grid}。
\begin{table}[!htbp]\centering\small
\caption{问题二各网格相对参考网格的最大差异}\label{tab:q2-grid}
\begin{tabular*}{0.96\textwidth}{@{\extracolsep{\fill}}ccc}
\toprule
网格区间数 & 温度最大差/$ {}^\circ\mathrm{C}$ & 含水率最大差/(kg/kg)\\\midrule
$60\times60$ & 0.0002 & 0.0004\\
$80\times80$ & 0.0001 & 0.0002\\
$100\times100$ & 0.0001 & 0.0001\\
$120\times120$ & 0.0001 & 0.0001\\
$140\times140$ & 0.0001 & 0.0001\\
$160\times160$ & 0.0000 & 0.0000\\
\bottomrule
\end{tabular*}
\par\smallskip{\footnotesize 注：参考网格为$160\times160$；差异按已有结果保留四位小数。}
\end{table}
$120\times120$与$160\times160$网格的两项差异均为$0.0001$，继续加密带来的变化较小，在当前输出精度下可以认为已经收敛；后续问题仍采用$120\times120$近壁加密网格。

\FloatBarrier
\subsection{问题三：烘干所需总时间计算}
问题三沿用问题二的模型，将计算延长至全域含水率达到干燥要求的时刻。
\subsubsection{环境条件的延拓}
附件1覆盖$0$至$14400\,\mathrm{s}$。设共有$N$组成对观测，取末$m$组的温度和水分浓度均值
\begin{equation}
\overline T_m=\frac1m\sum_{i=N-m+1}^{N}T_{\infty,i},\qquad
\overline C_m=\frac1m\sum_{i=N-m+1}^{N}C_{\infty,i}.
\end{equation}
环境边界写为
\begin{equation}\label{eq:tail}
(T_\infty(t),C_\infty(t))=
\begin{cases}
\operatorname{PCHIP}(\text{附件1数据})(t),&0\le t\le14400\,\mathrm{s},\\
(\overline T_m,\overline C_m),&t>14400\,\mathrm{s}.
\end{cases}
\end{equation}
正式方案取$m=20$，并在后文比较其他平均窗口。

\subsubsection{干燥终点与最湿位置}
\largepaperfigure[0.90]{q3_moisture_field.png}{3600秒时半子午区域内的含水率分布}{fig:q3-field}
\figref{fig:q3-field}展示圆柱体$r\ge0$、$z\ge0$的半子午截面，横轴为$r$，纵轴为$z$，白线表示等含水率线。在$3600\,\mathrm{s}$时，侧壁与端面附近的含水率较低，中心轴与中截面交点附近的含水率最高。题目要求药材各处含水率低于$0.15\,\mathrm{kg/kg}$，数值上以首次达到该阈值作为临界干燥时刻，定义
\begin{equation}\label{eq:event}
g(t)=\max_{(r,z)\in\Omega}C(r,z,t)-0.15,
\end{equation}
\begin{equation}
t_{\mathrm{dry}}
=\inf\left\{t\ge0:\max_{(r,z)\in\Omega}C(r,z,t)\le0.15\right\}.
\end{equation}
每次判断终止事件时，扫描全部二维节点，记录最大含水率及其位置。终点处的最湿节点由全域事件扫描确定，不以单个时刻的截面快照代替终点判定。

\subsubsection{烘干时间求解}
先按PCHIP边界积分至$14400\,\mathrm{s}$，再以末20点均值延拓边界，并接续该时刻的温湿场继续积分。BDF误差容限保持为$10^{-7}$和$10^{-9}$；当式\eqref{eq:event}首次沿下降方向达到零时，利用连续数值解定位终止时刻。

计算得到
\begin{equation}\label{eq:q3-time}
t_{\mathrm{dry},3}
=206767.5745\,\mathrm{s}
=57.4354\,\mathrm{h}.
\end{equation}
终止时全域最大含水率为$0.15\,\mathrm{kg/kg}$，位于几何中心。表面较早达到阈值；中心在$48\,\mathrm{h}$和$54\,\mathrm{h}$时的含水率仍分别为$0.1618$和$0.1538\,\mathrm{kg/kg}$，故总烘干时间取决于内部继续失水的过程。每隔$6\,\mathrm{h}$的含水率及终止时刻结果见\tabref{tab:q3-moisture}。
\begin{table}[!htbp]\centering\small
\caption{固定几何下中部截面的干基含水率，单位：$\mathrm{kg/kg}$}\label{tab:q3-moisture}
\renewcommand{\arraystretch}{1.15}
\begin{tabular*}{0.96\textwidth}{@{\extracolsep{\fill}}cccccc}
\toprule
 & \multicolumn{5}{c}{到药材轴线的距离/cm}\\
\cmidrule(lr){2-6}
时间/h & 0 & 0.5 & 1.0 & 1.5 & 2.0\\\midrule
6.0000 & 1.0168 & 0.9880 & 0.9014 & 0.7550 & 0.5329\\
12.0000 & 0.4560 & 0.4430 & 0.4029 & 0.3296 & 0.1642\\
18.0000 & 0.2987 & 0.2914 & 0.2686 & 0.2252 & 0.0844\\
24.0000 & 0.2377 & 0.2327 & 0.2167 & 0.1855 & 0.0664\\
30.0000 & 0.2058 & 0.2018 & 0.1891 & 0.1640 & 0.0598\\
36.0000 & 0.1858 & 0.1825 & 0.1718 & 0.1503 & 0.0566\\
42.0000 & 0.1720 & 0.1691 & 0.1597 & 0.1407 & 0.0548\\
48.0000 & 0.1618 & 0.1591 & 0.1506 & 0.1334 & 0.0537\\
54.0000 & 0.1538 & 0.1514 & 0.1436 & 0.1277 & 0.0530\\
57.4354 & 0.1500 & 0.1477 & 0.1402 & 0.1249 & 0.0526\\
\bottomrule
\end{tabular*}
\par\smallskip{\footnotesize 注：末行为烘干结束时刻，时间单位为h。}
\end{table}


\subsubsection{环境延拓敏感性分析}
在相同网格、密度、比热容、导热系数、水分扩散关系和误差容限下，比较不同末段平均窗口的烘干时间，见\figref{fig:tail-sensitivity}。以末20点方案为基准，末10点和末60点方案分别提前约$3.22\,\mathrm{min}$和延后约$2.42\,\mathrm{min}$；因此图中约$3.2\,\mathrm{min}$（约$0.09\%$）表示这组方案相对基准的最大绝对偏差，而非两端结果约$5.64\,\mathrm{min}$的跨度。单点延拓相对末20点方案提前约$16.22\,\mathrm{min}$。
\paperfigure{fig_q3_tail_sensitivity.png}{环境延拓点数对问题三烘干时间的影响}{fig:tail-sensitivity}
末20点覆盖约$19\,\mathrm{min}$，能够平均若干短时波动，同时反映数据截止前的状态。因此，后续采用
\begin{equation}
T_\infty=50.01765\unitC,\qquad
C_\infty=0.0499795\,\mathrm{kg/kg},\qquad t>14400\,\mathrm{s}.
\end{equation}
该方案假设烘房在数据截止后维持末段平均水平，敏感性结果仅适用于所比较的延拓情景。

\FloatBarrier
\subsection{问题四：径向收缩下的烘干模型}
问题四引入半径变化及题目附录4给定的密度、比热容、导热系数和水分扩散关系，环境延拓条件和干燥判据沿用问题三。
\subsubsection{径向收缩模型}
根据附件2的半径观测，以PCHIP构造$R(t)$。\figref{fig:radius}显示，半径在前期下降较快，后期趋于稳定。沿用长度不变、仅径向均匀收缩的假设，相对体积满足$V(t)/V_0=[R(t)/R_0]^2$。
\largepaperfigure{fig5_radius_shrinkage_dynamics.png}{药材半径、相对体积及横截面的变化}{fig:radius}

引入材料坐标
\begin{equation}\label{eq:material}
\xi=\frac r{R(t)},\qquad
\xi_i=1-\left(1-\frac i{N_r}\right)^2,\qquad
r_i(t)=R(t)\xi_i.
\end{equation}
节点在材料坐标中固定，实际位置随骨架移动。计算保留节点数与连接关系，按当前半径更新几何量。本问的密度、比热容、导热系数和水分扩散系数为
\begin{equation}\label{eq:q4-properties}
\begin{aligned}
\rho(C)&=760+90C,\\
c_p(C)&=1850+2150\frac C{C+1},\\
k(C)&=0.12+0.20\frac C{C+1},\\
D(T,C)&=4.2\times10^{-4}
\exp\left(-\frac{0.30}{C}\right)
\exp\left(-\frac{3850}{T+273.15}\right).
\end{aligned}
\end{equation}

\subsubsection{动态区域上的有限体积离散}
\paperfigure{moving_boundary_mesh_shrinkage.png}{径向收缩过程中网格的变化}{fig:moving-grid}
\figref{fig:moving-grid}给出相邻时刻的控制体。令$s_R(t)=R(t)/R_0$，由于轴向长度不变，控制体体积和界面面积满足
\begin{equation}\label{eq:geometry-scaling}
V_P(t)=s_R^2(t)V_P^0,\qquad
A^r(t)=s_R(t)A^{r,0},\qquad
A^z(t)=s_R^2(t)A^{z,0}.
\end{equation}
径向节点间距$d_r(t)=s_R(t)d_r^0$，轴向间距不变。因此，两方向的导通几何因子分别为
\begin{equation}
\frac{A^r(t)}{d_r(t)}=\frac{A^{r,0}}{d_r^0},\qquad
\frac{A^z(t)}{d_z}=s_R^2(t)\frac{A^{z,0}}{d_z}.
\end{equation}
将动态几何量代入统一离散格式：
\begin{equation}\label{eq:moving-fvm}
b_{u,P}V_P(t)\frac{\mathrm du_P}{\mathrm dt}
=\sum_{N\in\mathcal N_P}
\frac{a_{u,f}A_f(t)}{d_{PN}(t)}(u_N-u_P)
+\sum_b\beta A_{b,P}(t)(u_\infty-u_P).
\end{equation}
时间导数在固定材料坐标下计算。每次计算方程右端和雅可比矩阵时，按当前半径更新体积、面积及间距，按当前温湿状态更新密度、比热容、导热系数和水分扩散系数。该形式以温度和干基含水率作为材料场变量，描述给定收缩轨迹下的局部扩散过程。

\subsubsection{烘干时间求解与结果分析}
计算取$N_r=N_z=120$、$\lambda=2$，环境采用末20点均值延拓，终止事件为
\begin{equation}
g_4(t)=\max_P C_P(t)-0.15.
\end{equation}
同时更新两个物理场与半径，每隔$60\,\mathrm{s}$保存结果。固定输出位置$r^\ast$仅在$r^\ast\le R(t)$时具有药材内部的物理意义；域外位置不输出数值，另行记录当前表面含水率。

计算得到
\begin{equation}\label{eq:q4-time}
t_{\mathrm{dry},4}
=183809.7448\,\mathrm{s}
=51.0583\,\mathrm{h}.
\end{equation}
终止时半径为$1.2000\,\mathrm{cm}$，中心含水率为$0.1500\,\mathrm{kg/kg}$，表面为$0.0526\,\mathrm{kg/kg}$。题目指定时刻、位置的结果见\tabref{tab:q4-moisture}。
\begin{table}[!htbp]\centering\small
\caption{径向收缩下中部截面的干基含水率，单位：$\mathrm{kg/kg}$}\label{tab:q4-moisture}
\renewcommand{\arraystretch}{1.15}
\begin{tabular*}{0.96\textwidth}{@{\extracolsep{\fill}}cccccccc}
\toprule
 & \multicolumn{6}{c}{到药材轴线的距离/cm} & \\
\cmidrule(lr){2-7}
时间/h & 0 & 0.5 & 1.0 & 1.5 & 2.0 & 药材表面 & 半径/cm\\\midrule
6.0000 & 1.7186 & 1.5368 & 1.0220 & --- & --- & 0.4206 & 1.3740\\
12.0000 & 0.7370 & 0.6541 & 0.4080 & --- & --- & 0.1668 & 1.2480\\
18.0000 & 0.4084 & 0.3684 & 0.2396 & --- & --- & 0.0891 & 1.2140\\
24.0000 & 0.2849 & 0.2609 & 0.1789 & --- & --- & 0.0673 & 1.2040\\
30.0000 & 0.2262 & 0.2093 & 0.1492 & --- & --- & 0.0595 & 1.2010\\
36.0000 & 0.1927 & 0.1796 & 0.1316 & --- & --- & 0.0559 & 1.2000\\
42.0000 & 0.1712 & 0.1603 & 0.1199 & --- & --- & 0.0541 & 1.2000\\
48.0000 & 0.1561 & 0.1467 & 0.1115 & --- & --- & 0.0530 & 1.2000\\
51.0583 & 0.1500 & 0.1413 & 0.1081 & --- & --- & 0.0526 & 1.2000\\
\bottomrule
\end{tabular*}
\par\smallskip{\footnotesize 注：含水率单位为kg/kg，半径单位为cm；“---”表示该位置已在药材外部，末行为烘干结束时刻。}
\end{table}


\largepaperfigure[1.18]{q4_combined_pizza.png}{径向收缩下的含水率与温度变化（左：含水率；右：温度）}{fig:q4-fields}
\figref{fig:q4-fields}沿用\figref{fig:q1-fields}的展开方式，外缘随时间向内移动，颜色表示含水率或温度，两幅图使用各自的时间刻度。药材半径从$2.0\,\mathrm{cm}$逐渐减小，$36\,\mathrm{h}$后基本稳定于$1.2\,\mathrm{cm}$。$6\,\mathrm{h}$时半径已为$1.374\,\mathrm{cm}$，因此$1.5\,\mathrm{cm}$和$2.0\,\mathrm{cm}$位置已位于药材外部；中心仍是最后达到阈值的位置。

与问题三相比，问题四的总时间减少
\begin{equation}
\Delta t=t_{\mathrm{dry},3}-t_{\mathrm{dry},4}
=22957.8297\,\mathrm{s}=6.3772\,\mathrm{h},
\end{equation}
\begin{equation}
\frac{\Delta t}{t_{\mathrm{dry},3}}\times100\%=11.10\%.
\end{equation}
两问同时改变半径以及密度、比热容、导热系数和水分扩散关系，因此$11.10\%$是二者共同作用下的时间差。若要单独量化收缩影响，应在问题四的上述参数关系下增加固定半径对照。
\endgroup

\section{模型评价与改进}
\subsection{模型的优点}
\begin{enumerate}[label=\arabic*.]
\item 有限体积法使公共界面两侧通量大小相等、方向相反，内部通量在全域求和时相互抵消，便于检查离散方程的收支一致性。
\item 四个问题依次加入温湿耦合、全域终止事件和动态几何更新，共用网格与边界处理方式，计算框架具有较强的复用性。
\item 二维轴对称模型同时描述侧壁与端面的热量和水分交换，并通过全域扫描确定最湿位置。
\item 近壁加密、网格收敛性比较及环境延拓敏感性分析，分别检验了高梯度区域离散、网格规模和长时边界选择对结果的影响。

\end{enumerate}

\subsection{模型的不足与改进方向}
\begin{enumerate}[label=\arabic*.]
\item 温度方程未计入水分蒸发的耗热效应，后续可增加潜热项及相应的热质耦合关系。
\item 模型未反映药材纤维与孔隙结构的方向差异，后续可引入方向相关的导热系数和水分扩散系数。
\item 数据截止后的环境由末段均值给定，所考察方案间的差异较小不能替代长时间实测边界数据。
\item 问题四按给定半径轨迹更新几何，后续应完善干物质密度与骨架运动的相容关系，并在相同参数关系下增加固定半径对照，以单独量化收缩影响。
\end{enumerate}

\clearpage
\section*{AI工具使用说明}
本参赛队在竞赛过程中使用了 AI 工具，主要用于数学模型与数值程序的辅助审查、代码审计及数值验证方案检查，详细使用情况见支撑材料。
% ============================================================
% 参考文献
% ============================================================
% 正文中使用 \cite{文献键}；仅被引用的条目会出现在最终参考文献表中。
\bibliographystyle{gbt7714-numerical}
\bibliography{references}


\end{document}
